\documentclass[11pt,a4paper]{article}
\usepackage[T1]{fontenc}
\usepackage[utf8]{inputenc}
\usepackage[brazil]{babel}
\usepackage[margin=2.2cm]{geometry}
\usepackage{booktabs,longtable,array}
\usepackage[hidelinks]{hyperref}
\usepackage{xurl}

\title{Relatório de Limpeza de Dados\\\large Análise de Dados de Enchentes do INCT SiMAI}
\author{Henrique Nunes (\textit{completar com os demais integrantes})\\Coleta, Preparação e Análise de Dados --- PUCRS 2026/2}
\date{9 de novembro de 2026}

\begin{document}
\maketitle

\section{Objetivo e reprodução}
Este relatório descreve as inconsistências encontradas nos dados usados pelo dashboard, onde aparecem, por que são um problema e como foram tratadas. Os números de pessoas e prédios são \textbf{estimativas de uma simulação}, não previsões nem contagens oficiais.

Para reproduzir, execute na raiz do repositório:
\begin{verbatim}
python scripts/coleta_ana.py
python scripts/limpa_inct.py
python scripts/processa_ana.py
\end{verbatim}
Cada alteração nos prédios afetados está registrada em \texttt{dados/processados/inct\_inconsistencias.csv}. Os valores originais continuam no arquivo limpo, nas colunas com sufixo \texttt{\_original}.

\section{JSON do INCT}
O arquivo tem 71 níveis (3,0 a 10,0~m) e 94 bairros, em formato longo de 6.674 linhas. Não encontramos níveis com número diferente de bairros, nomes duplicados ou com espaços sobrando, valores ausentes, \texttt{affected\_buildings} maior que \texttt{total\_buildings} ou \texttt{total\_buildings}/\texttt{population} variando entre níveis.

\subsection{Prédios afetados que caem quando o rio sobe}
\textbf{Problema.} Se o nível sobe, a área alagada não diminui, então \texttt{affected\_buildings} não pode cair. Há 12 quedas em 5 bairros.

\textbf{Tratamento: corrigir.} Para cada bairro, mantemos a maior sequência não decrescente de valores. Os valores fora dela são considerados erro e interpolados linearmente entre os vizinhos válidos; nas pontas, repetimos o valor válido mais próximo. Essa regra foi escolhida porque um teste simples de ``nunca descer'' trataria os valores altos errados como corretos e apagaria os valores seguintes. Foram corrigidos 17 valores.

\begin{longtable}{l c r r p{6cm}}
\toprule
Bairro & Nível (m) & Original & Corrigido & Observação\\
\midrule
\endhead
Espírito Santo & 3,0--3,3 & 637 & 25 & Bloco alto antes de cair para 25 em 3,5~m\\
Espírito Santo & 3,4 & 633 & 25 & Mesmo bloco\\
Espírito Santo & 4,7 & 2.101 & 46 & Salto isolado; 46 em 4,6~m e 47 em 4,8~m\\
Centro Histórico & 9,7 & 0 & 2.077 & Zerado entre 2.050 e 2.130\\
Centro Histórico & 9,8 & 0 & 2.103 & Zerado entre 2.050 e 2.130\\
Rubem Berta & 8,8 & 0 & 268 & Zerado entre 252 e 285\\
Lageado & 8,1 & 0 & 109 & Zerado entre 108 e 110\\
Arquipélago & 3,0; 3,5; 4,0; 4,5; 5,0 & 4.685 & 4.680 & Ver seção~\ref{sec:rodadas}\\
Arquipélago & 4,1 & 4.674 & 4.680 & Oscilação pequena\\
Arquipélago & 5,3 & 4.688 & 4.680 & Ver seção~\ref{sec:rodadas}\\
\bottomrule
\end{longtable}

\textbf{Incerteza.} No Espírito Santo, os níveis de 3,0 a 3,4~m seguem a maioria da série (25 prédios), mas não há como confirmar qual dos dois patamares é o correto. Mantivemos a correção e pedimos a verificação ao INCT.

\subsection{\texttt{percentage} arredondado para cima}
\textbf{Problema.} O campo é inteiro e sempre arredondado para cima: a diferença para o valor exato chega a 1 ponto percentual, e 2.638 das 6.674 linhas foram arredondadas. Em Santa Tereza a 3,2~m, 1 prédio em 15.906 vira 1\%.

\textbf{Tratamento: corrigir.} Recalculamos \texttt{percentage} como $100 \times$ \texttt{affected\_buildings} / \texttt{total\_buildings}, com duas casas decimais.

\textbf{Efeito de ponto flutuante.} Em Boa Vista do Sul a 7,2 e 7,3~m, 686 prédios em 2.450 dão exatamente 28\%, mas o arquivo traz 29\%. O resultado é compatível com um arredondamento para cima aplicado a um valor como 28,000000000000004. O recálculo também resolve esse caso.

\subsection{\texttt{affected\_population} herda o arredondamento}
\textbf{Problema.} O campo segue a fórmula do enunciado, mas usa o percentual arredondado, e também é arredondado para cima. Santa Tereza a 3,2~m tem 314 pessoas afetadas por 1 prédio. No total da cidade, nos 55 níveis sem correção de prédios, o excesso vai de 1.326 a 3.963 pessoas (mediana de 2.744).

\textbf{Tratamento: corrigir.} Recalculamos com o percentual exato: \texttt{population} $\times$ \texttt{affected\_buildings} / \texttt{total\_buildings}, arredondado ao inteiro mais próximo. Santa Tereza passa a 2 pessoas. A 5,4~m, o total da cidade vai de 141.108 para 139.179 pessoas.

\subsection{Ordem dos campos e indícios de rodadas diferentes}\label{sec:rodadas}
\textbf{Problema.} Em 6 níveis (3,0; 3,5; 4,0; 4,5; 5,0 e 5,3~m), os campos aparecem em outra ordem. A 5,3~m a ordem dos bairros também muda. A leitura por nome não é afetada, mas o Arquipélago oscila exatamente nesses níveis (4.685 ou 4.688, contra 4.680 nos demais). Isso sugere que esses níveis vêm de outra rodada ou versão da simulação.

\textbf{Tratamento: manter e avisar.} A ordem é irrelevante para o formato longo. Não sabemos se os demais valores desses níveis diferem, então essa questão faz parte da devolutiva ao INCT.

\subsection{Limitações mantidas com aviso}
\begin{itemize}
  \item A população é distribuída igualmente entre os prédios, inclusive os não residenciais.
  \item A simulação depende apenas do nível do rio. Diques, comportas e casas de bomba não aparecem nos dados.
  \item Bairros como Arquipélago (99,83\%) e Farrapos (79,44\%) já aparecem muito afetados a 3,0~m. Mantivemos esses valores por serem coerentes com a série, com aviso no dashboard.
\end{itemize}

\section{Telemetria da ANA}
Série de 15 minutos das estações Cais Mauá C6 (\texttt{87450004}) e Usina do Gasômetro (\texttt{87444000}), coletada em 28/10/2026.

\begin{longtable}{p{4.2cm} p{5.6cm} p{5.6cm}}
\toprule
Inconsistência & Onde & Tratamento\\
\midrule
\endhead
Picos falsos do sensor & 123 leituras do Cais Mauá (2015: 8; 2017: 39; 2019: 57; 2025: 2; 2026: 17), incluindo o suposto 3,90~m de 2017 & Descartar leituras a mais de 20~cm da mediana de 9 leituras. O pico oficial de 5,35~m é verificado após a limpeza\\
Sensor travado & 69 dias entre 29/10/2015 e 20/01/2016, com patamares de 268~cm e 98~cm & Descartar dias cuja janela de 15 dias varia até 2~cm\\
Registros sem nível & Tags \texttt{<Nivel/>} vazias & Ignorar na leitura do XML\\
Falta de leitura do Cais Mauá na enchente & Última leitura em 02/05/2024 às 23h15 (369~cm); retorno em 11/06/2025 às 23h45 & Usar o Gasômetro entre 03/05/2024 e 11/06/2025, com a coluna \texttt{estacao}\\
Lacuna entre as estações & 18 horas sem leitura entre 02/05 e 03/05/2024 & Interpolar linearmente e marcar com \texttt{origem = interpolado}\\
Referenciais diferentes & O Gasômetro lê entre 41 e 80~cm acima do Cais Mauá nos períodos em comum; nota ANA/SGB cita correção de 1,18~m & Manter sem ajuste e avisar. O pico servido bate com os 5,35~m oficiais\\
Cotas de inundação diferentes & 3,00~m no Cais Mauá; 3,60~m no Gasômetro & Comparar cada dia com a cota da estação em uso\\
Lacunas maiores & Por exemplo, 85 dias até 14/10/2015 e 70 dias até 05/07/2016 & Manter em branco; ressaltar que os anos não têm a mesma cobertura\\
\bottomrule
\end{longtable}

\section{Ligação entre as bases}
Cada nível da ANA é arredondado ao décimo mais próximo, com meio para cima (5,35 vira 5,4~m), e consultado no INCT. O resultado é uma estimativa do que aquele nível representa na simulação, e não uma contagem de quem foi atingido em cada hora. Abaixo de 3,0~m não há simulação e o campo fica vazio.

\section{Pendências}
\begin{itemize}
  \item Confirmar com o INCT se a simulação usa o referencial do Cais Mauá ou do Gasômetro.
  \item Enviar ao INCT a devolutiva: quedas de prédios, arredondamento e os 6 níveis com estrutura diferente.
\end{itemize}

\end{document}
